%document class
\documentclass[10pt,oneside]{book}
%%%% Page Info + Commands %%%%%{

%packages
\usepackage{pgfplots}
\pgfplotsset{compat=1.18}
\usepackage{amsmath}
\usepackage{amsfonts}
\usepackage{charter}
\usepackage{amsthm,letltxmacro}
\usepackage{amssymb}
\usepackage{fancyhdr}
\usepackage{float}
\usepackage[most]{tcolorbox}
\usepackage{enumerate}
\usepackage{enumitem}
\usepackage[dvipsnames]{xcolor}
\usepackage{changepage}
\usepackage{mdframed}
\usepackage{graphicx}

% SMALL PAGE COMMAND
\newcommand{\smallpage}{  \setlength \oddsidemargin{.25in}
            \setlength \textwidth{5.5in}
            \setlength \topmargin{0in}
            \setlength \textheight{9in}}


% Commands for sets of numbers
\newcommand{\Z}{\mathbb{Z}}\newcommand{\R}{\mathbb{R}}\newcommand{\N}{\mathbb{N}}\newcommand{\Q}{\mathbb{Q}}\newcommand{\C}{\mathbb{C}}
\newcommand{\real}[1]{\text{Re}\left(#1\right)}
\newcommand{\im}[1]{\text{Im}\left(#1\right)}

% gordie spacing and boxing commands
\newcommand{\gspc}{\vspace{0.13cm}} % 0.5 space
\newcommand{\gline}[0]{\gspc\\} % 1.5 space
\newcommand{\gbline}{\gspc\gspc\\} % double space
\newcommand{\gbox}[1]{\fbox{\parbox{\linewidth}{#1}}} % multiline box command
\newcommand{\gmod}[1]{\,(\text{mod}\,#1)}


\newcommand{\gimage}[2]{
\begin{figure}[H]
    \centering
    \includegraphics[width=#2]{#1}
\end{figure}
}

\newcommand{\centerit}[1]{{\begin{center} #1 \end{center}}}
\newcommand{\ul}[1]{\underline{#1}}


% COLOR BOX
\newmdenv[backgroundcolor=gray!8,
  linewidth=0pt, innerleftmargin=0pt, innerrightmargin=0pt,
  skipabove=\baselineskip, skipbelow=\baselineskip
]{coloredadjust}

% PROOF
\LetLtxMacro\oldproof\proof
\let\endoldproof\endproof
\renewenvironment{proof}[1][\proofname]
  {\vspace{0.2cm}\begin{adjustwidth}{2em}{2em}
   \oldproof[#1]}
  {\endoldproof
\end{adjustwidth}}

% PROBLEM
\newenvironment{problem}[1]
  {\vspace{-0.1cm}\begin{coloredadjust}\begin{adjustwidth}{2em}{2em}
   \textit{Problem. }#1}
  {\endoldproof
\end{adjustwidth}\end{coloredadjust}\gspc}

%% DEFINITION
\newtcbtheorem[]{definition}{Definition}{enhanced,
  colback=blue!2, colframe=blue!50!black, coltitle=blue!70!black,
  fonttitle=\bfseries,
  boxrule=0pt, toprule=2pt, bottomrule=1pt, leftrule=1pt, rightrule=1pt, arc=1pt, outer arc=1pt,
  attach title to upper, title={#1},
}{dfn}

%% CRITERION
\newtcbtheorem[]{criterion}{Criterion}{enhanced,
  colback=magenta!3, colframe=magenta!60!black, coltitle=magenta!20!black,
  fonttitle=\bfseries,
  boxrule=4pt, toprule=2pt, bottomrule=1pt, leftrule=1pt, rightrule=1pt, arc=1pt, outer arc=1pt,
  attach title to upper, title={#1},
}{ctn}

%% COROLLARY
\newtcbtheorem[]{corollary}{Corollary}{enhanced,
  colback=orange!3, colframe=orange!80!black, coltitle=orange!50!black,
  fonttitle=\bfseries,
  boxrule=4pt, toprule=1pt, bottomrule=1pt, leftrule=1pt, rightrule=1pt, arc=5pt, outer arc=5pt,
  attach title to upper, title={#1},
}{ctn}

%% THEOREM
\newtcbtheorem[]{theorem}{Theorem}{enhanced,
  colback=green!2, colframe=green!40!black, coltitle=green!20!black,
  fonttitle=\bfseries,
  boxrule=4pt, toprule=3pt, bottomrule=1pt, leftrule=1pt, rightrule=1pt,arc=5pt, outer arc=5pt,
  attach title to upper, title={#1},
}{thm}

%%%% Page Setup %%%%%%%
\smallpage\sloppy
\pagestyle{fancy}
\lhead{\large{\textbf{Feb 9th{}-{}Notes}}}
%\rfoot{\thepage/\pageref{LastPage} }
\setlength{\headheight}{10pt}

\begin{document}

\newcommand{\cg}[1]{\color{OliveGreen}#1\color{white}}
\newcommand{\cb}[1]{\color{BlueViolet}#1\color{white}}

\subsection*{Computability}
It is possible that there are some programs that are impossible or \textbf{not computable}. 

In turn, this makes us think about some general questions:
\begin{itemize}
  \item What kinds of things are \textbf{computable}? That is, what is possible fo a computer to do? What is impossible?
  \item What does ``computable'' mean?
  \item What is a computer?
\end{itemize}

\subsection*{Deterministic Finite Automata}
Consider a trffic light, which is like a simple computer. 

We have three states: red, yellow, and green. 
\begin{itemize}
  \item If we are red, we can only move to green
  \item If we are green, we can only move to yellow
  \item If we are yellow, we can only move to red. 
\end{itemize}
\begin{definition}{Deterministic Finite Automata}{}\gline{}
  A Deterministic finite automaton consist of:
  \begin{itemize}[itemsep=0pt, parsep=0pt]
    \item states
    \item an alphabet
    \item rules for transitions from each tate ofr each letter in the alphabet
  \end{itemize}
  And also:
  \begin{itemize}[itemsep=0pt, parsep=0pt]
    \item One state chosen to be the start state
    \item States chosen to be the accept states
  \end{itemize}
\end{definition}
We diagram the states like this:
\gimage{image.png}{7cm}
\begin{definition}{Deterministic Finite Automata}{}\gline{}
  A deterministic finite automaton consists of:
  \begin{itemize}[itemsep=1pt, parsep=1pt]
    \item States
    \item An alphabet
    \item Rules for transition between each state
  \end{itemize}
  We can write this more mathematically:\gline{}
  A (deterministic) \textbf{finite automaton} is a 5-tuple ($Q, \sum, \delta, q_0, F$). 
  \begin{itemize}[itemsep=1pt, parsep=1pt]
    \item $Q$ is a finite set of states
    \item $\sum$ a finite set of input symbols
    \item $\delta$ the transition function
    \item $q_0\in Q$ the starting state
    \item $F = \{q_i, \cdots\}$ Set of accepting states. 
  \end{itemize}
\end{definition}
\begin{definition}{Regular Languages}{}\gline{}
  A language is caled a \textbf{regular language} if some finite automaton recognizes it.
  
\end{definition}
\section*{Feb 12th}
\subsection*{Combinations of Languages}
Today we are studing how to combine different regular languages together. A lot of these operations should be familiar if you took Discrete mathematics. The ones I think we will be using are below:\gline{}
Union: $A\cap B = \{x, y\mid x\in A\text{ and } y\in B\}$\\
Concatenation: $A\circ B = \{xy \mid x\in A, y\in B\}$\\
Star: $A^*=\{x_1x_2x_3\cdots x_k\mid k\ge 0\text{ and each }x_1\in A\}$
\subsection*{Concatenation}
Let $A = \{1, 11, 111\}$ and $B=\{0,00\}$. Then:
\begin{align*}
  A\circ B &= \{10, 100, 110, 1100, 1110, 11100\}\\
  B\circ A &= \{01, 011, 0111, 001, 0011, 00111\}
\end{align*}
Let $A\{0, 1, 01\}$ And $B = \{1, 01, 000\}$. Then:
\begin{align*}
  A \circ B &= \{01, 001, 0000, 11, 101, 1000, 011, 0101, 01000\}\\
  B\circ A &= \{10, 11, 101, 010, 011, 0101, 0000, 0001, 00001\}
\end{align*}
\subsection*{Star}
A* = $\{1, 11, 111\}.$ Then: 
\begin{align*}
  A^*=\{1, 11, 111, \cdots\}
\end{align*}
Let $B^* = \{0, 11\}$. Then:
\begin{align*}
 B^* = \{0, 11, 011, 110, 1111, 11011, \cdots\}
\end{align*}
Let $C^* = \{0, 1\}$. Then:
\begin{align*}
  C^* = \{0, 1, 01, 10, 010, \cdots\}
\end{align*}
For C* in particular, we typically call it ${\sum}^*=\{0, 1\}^*$.\gline{}
A simplification that we can use is that for any $A^*$, we have that:
\begin{align*}
  A^* &= A\cup A \cup A\circ A \cup A \circ A \circ A \cup \ldots\\
  &= \bigcup_{1=0}^\infty A^i
\end{align*}
Now, we're going to work on drawing some DFAs.\gline{}
Let $A$ be the set of binary string that contains at least two 0's\gline{}
Let $B$ be the set of binary string that end in 01.
\begin{enumerate}
  \item Draw a diagram for a DFA that recognizes $A\circ B$. 
  \gimage{image2.png}{8cm}
  \item Draw a driagram for a DFA that recognizes $A\cup B$.
  \gimage{image3.png}{8cm}
\end{enumerate}
\end{document}